\clearpage
\section*{四、实验心得与思考题}

\subsection*{1. 实验心得}

通过这次实验，我得到了以下收获与心得：

\begin{enumerate}
\item 我学会了密立根油滴仪“平衡—提升—0V”的操作流程，体会到让一颗油滴在视场中往复运动并保持清晰需要耐心。
\item 选择油滴是关键。本次油滴带电量较多（26～140 个电子），差值类方法的误差随之增大，这一点需要改进。
\item 用三种方法处理同一组数据后，我发现平均值法和公约数法结果稳定，而多重差值法在带电量较大时误差明显放大，说明数据处理方法要与数据特点相匹配。

\end{enumerate}
\subsection*{2. 思考题}

按要求从五道题中任选三题，我选择第 2、3、4 题。

\subsubsection*{1. 为何选油而不用水？}

\begin{enumerate}[leftmargin=3.6em]
\item 水易蒸发，微米级水滴在测量过程中质量会明显减小，平衡电压和下落速度随之漂移，无法稳定测量。实验用的钟油蒸气压很低，挥发极慢，质量在测量期间基本不变。
\item 钟油的密度已知且稳定，表面张力使油滴保持球形，满足斯托克斯定律的假设，便于由下落速度求半径。

\end{enumerate}
\subsubsection*{2. 如何从一堆 $q$ 的测量值看出电荷量子化？}

\begin{enumerate}[leftmargin=3.6em]
\item 把各油滴的 $q$ 从小到大排列，若电荷是量子化的，$q$ 不会连续分布，而是集中在某个最小值的整数倍附近，相邻差值也都接近这个最小值的整数倍。
\item 也可以用 $q/e_0$ 检验：若各值都接近整数，就说明 $q$ 是 $e_0$ 的整数倍。作 $q$ 与 $n$ 的关系图时，各点应落在一条过原点的直线上，斜率即为 $e$。
\item 这种判断要求测量误差明显小于一个电子的电荷，因此选用带电少的油滴时，$q$ 的分立特征更容易直接看出。

\end{enumerate}
\subsubsection*{3. 浮力能忽略吗？}

\begin{enumerate}[leftmargin=3.6em]
\item 考虑空气浮力时，重力应换成 $\frac43\pi r^3(\rho-\rho_{\text{空}})g$。空气密度约 $1.2\,\mathrm{kg/m^3}$，只有钟油密度 $981\,\mathrm{kg/m^3}$ 的约 0.12\%。
\item 代入推导可知 $q\propto(\rho-\rho_{\text{空}})^{-1/2}$，忽略浮力只使 $q$ 偏小约 0.06\%。
\item 这比本次同一油滴两次测量约 1\%～8\% 的差别小得多，因此浮力可以忽略。
\end{enumerate}
